\section{アプローチ08（縦方向乱立対策およびWeb環境パフォーマンス最適化手法）}

\subsection{手法の概要}
アプローチ07の重複排除をすり抜けて発生する「縦（Y軸）方向に乱立する小さなバウンディングボックス」への対策と、持続的な実行におけるメモリ蓄積・処理遅延を抑制するパフォーマンス最適化を導入した。

重複対策を通過した全人物のバウンディングボックスに対し、二次元平面上のY方向およびX方向の重なり具合を検証する。X座標の重なりが設定閾値を超え、かつY方向に連続して位置するバウンディングボックス群を同一人物の誤断片化と判定し、それらをすべて包摂する単一のバウンディングボックスへ結合・統一する処理を適用した。

また、Web環境（TensorFlow.js / WebGL / WebGPU）における実行パフォーマンスとOOM（Out of Memory）対策を徹底するため、以下の環境最適化を実施した。

\begin{itemize}
    \item \textbf{WebGPUおよびWebGL環境の最適化}：WebGPUが利用可能な環境では優先的に適用し、WebGL環境下においては 16ビット浮動小数点推論（\texttt{WEBGL\_RENDER\_FLOAT32\_CAPABLE} を \texttt{false}、\texttt{WEBGL\_FORCE\_F16\_TEXTURES} を \texttt{true} に設定）を採用して計算負荷とメモリ占有量を削減した。
    \item \textbf{テクスチャキャッシュの即時解放}：入力画像の内容およびサイズが毎フレーム変動することに伴う不必要なキャッシュの蓄積を完全に防ぐため、\texttt{WEBGL\_DELETE\_TEXTURE\_THRESHOLD} を \texttt{0} に設定し、使い捨てテクスチャの即時破棄を強制した。
\end{itemize}

\subsection{結果およびアプローチ07との比較・課題}

\subsubsection{アプローチ07との比較}
縦方向に細かく乱立していた不正なバウンディングボックスが適正に一元化され、ポーズ検出結果を保持した正確な人物領域のみが残るようになった。さらに、16ビット推論への移行とテクスチャ即時解放の導入により、長時間の連続実行においてもメモリリーク（OOM）を起こすことなく安定した高フレームレートを維持できる環境が整った。

\subsubsection{判明した課題}
\begin{enumerate}
    \item \textbf{人物の単体検出にとどまる構造的限界}：
    バウンディングボックスの統合および最適化により個体の特定精度は大きく高まったが、依然として人物が「待機列内のどのような並び（ライン）を形成しているか」や「誰と誰が向き合ってグループを構成しているか」という集合レベルの構造判定には至っていなかった。
    \item \textbf{幾何学的情報・向き情報の未活用}：
    ポーズ検出によって得られた関節座標情報をバウンディングボックスの補正のみに使用しており、グループ判定そのものは幾何学的な位置関係やベクトルを活用できていない課題が残された。
\end{enumerate}
